\bookchapter{結語：保持好奇，保持清醒}

恭喜你讀到了這裡。

這本書從手機硬體拆到軟體生態，從日常設定挖到開發者模式，從安全防護聊到數字遺產。如果它幫你解決了一個實際問題——哪怕只是學會了關掉一個不該開的許可權、養成了一次備份的習慣——那這些字就沒白寫。

但比任何具體技巧都重要的，是兩件事：

\begin{itemize}
  \item \textbf{保持好奇}：手機每年都在變，AI在變，摺疊屏在變，衛星通訊在變。今天學的操作步驟明天可能就過時了，但"想知道為什麼"的好奇心永遠不會過時。多問一句"這個許可權是幹什麼的""這個資料存在了哪裡""這個功能會不會洩露隱私"，比記住一百個操作步驟都有用。
  \item \textbf{保持清醒}：手機是工具，不是主人。它設計出來就是為了吸引你的注意力、佔據你的時間、引導你的消費。用它是你的選擇，被它用是你的失控。定期抬頭看看螢幕時間，定期想想"我是在用手機做事，還是在被手機消磨"，這不是矯情，這是數字時代的生存本能。
\end{itemize}

技術在變，但人的需求沒變——我們需要安全感、需要連線、需要被尊重。手機只是一個載體。真正的“手機課”不在書裡，在你每天解鎖螢幕的那一瞬間。

\begin{tcbpractice}{最後一個練習}
合上這本書，拿起你的手機，做一件事：\textbf{用30秒想想，你的手機裡有哪三樣東西是“丟了會真正心疼”的？}然後，確保它們被備份了。這本書就值了。

\end{tcbpractice}
